\section{Introducción}
\label{sec:introduccion}

La memoria documenta el desarrollo de la práctica final de la
asignatura \textit{Aprendizaje en Inteligencia Artificial} del Grado en
Ingeniería en Inteligencia Artificial de la Universidad Rey Juan Carlos.
El proyecto, denominado internamente \textbf{PFINAL\_AP-IA}, constituye
la finalización e integración de las seis prácticas anteriores
(P1--P5) bajo un único sistema multiagente conversacional dedicado a la
gestión financiera personal.

El sistema toma la forma de un asistente financiero capaz de mantener
un diálogo natural con el usuario, registrar transacciones a partir de
texto o imágenes (OCR), detectar anomalías de gasto, generar
predicciones temporales, autenticar mediante biometría facial y
explicar sus respuestas mediante visualizaciones dinámicas. La
arquitectura se apoya en \textbf{LangGraph}~\cite{langgraph} para la
orquestación de agentes, \textbf{FastAPI}~\cite{fastapi} para exponer
cada módulo como microservicio REST, \textbf{Next.js}~\cite{nextjs}
en el frontend y \textbf{Langfuse}~\cite{langfuse} para la
monitorización de trazas, llamadas a herramientas, latencias y
consumo de tokens.

El documento se estructura en cinco capítulos. El capítulo de
\textit{Contexto} sitúa la práctica dentro del recorrido de las prácticas anteriores y
recoge los requisitos funcionales. El \textit{Estado del arte} repasa
los frameworks y técnicas que sustentan las decisiones de diseño. La
\textit{Propuesta} detalla la arquitectura implementada, los dos roles
diferenciados (usuario y administrador), la memoria conversacional, la
distribución de los módulos P1--P5 como microservicios, la evolución
incremental sobre P2, P3 y P5, las metodologías de prompt engineering
empleadas y la estrategia de pruebas unitarias y de estrés. Las
\textit{Conclusiones} resumen los resultados cuantitativos obtenidos y
las líneas de trabajo futuro.
